\SourceLeaf{392}%
\[
\begin{aligned}
\mathrm{V.}\quad&2aa\sin2(A''-A')+2a'a'\sin2(A-A'')+2a''a''\sin2(A'-A)\\
&\quad+t[\sin2(A''-A')+\sin2(A-A'')+\sin2(A'-A)]\\
&\quad-4ar\cos(A-\varphi)\sin2(A''-A')\\
&\quad-4a'r\cos(A'-\varphi)\sin2(A-A'')\\
&\quad-4a''r\cos(A''-\varphi)\sin2(A'-A)=0
\end{aligned}
\]
\par
Das Aggregat, worin hier \(t\) multiplicirt erscheint, kann auch durch
\[
4\sin(A''-A')\sin(A''-A)\sin(A'-A)
\]
ausgedrückt werden.

Behandelt man auf eine ähnliche Art die Gleichungen I, II, IV, so bekommt man eine ähnliche Gleichung VI, die sich von V nur durch die Vertauschung der Buchstaben \(a'',A''\) gegen \(a''',A'''\) unterscheidet. Eliminirt man aus den beiden Gleichungen V und VI die Grösse \(t\), so sieht man leicht, dass daraus eine Gleichung von der Form
\[
\mathrm{VII.}\qquad B+Cr\cos\varphi+Dr\sin\varphi=0
\]
hervorgehen wird, wo \(B,C,D\) bekannte Grössen bedeuten. Man kann ihre Werthe leicht darstellen, wir werden indess bald zeigen, wie man dieser Entwickelung überhoben sein kann. Aus der Gleichung VII ist klar, dass der Mittelpunkt jeder die vier Seiten unsers Vierecks berührenden Ellipse in einer geraden Linie liegt, welche gegen die Abscissen-Linie unter einem Winkel, dessen Tangente \(=-\frac{C}{D}\), geneigt ist, und dass der Durchschnitts-Punkt die Abscisse \(-\frac{B}{C}\) hat. Die Lage dieser geraden Linie kann man aber viel leichter durch folgende Betrachtungen bestimmen. Eine Diagonale des Vierecks kann als eine verschwindende, die Seiten des Vierecks berührende Ellipse betrachtet werden, deren Mittelpunkt dann offenbar in der Mitte der Diagonale liegt. Hieraus folgt leicht, dass die obige gerade Linie, welche der geometrische Ort der Mittelpunkte aller die vier Seiten des Vierecks berührenden Ellipsen ist, keine andere sein könne, als die, welche die Halbirungspunkte der beiden Diagonalen verbindet, und welche demnach leicht gefunden werden kann. Hierüber füge ich noch zwei Bemerkungen hinzu:

1) Fielen beide Halbirungs-Punkte in Einen zusammen (in welchem Falle das Viereck ein Parallelogramm sein wird), so fällt freilich diese Bestimmung der
